\section{Problem Statement}

We consider a stream of identical boxes arriving on a belt.
Each box has size \(x \times y\), where \(x\) is parallel to the
belt direction and \(y\) is orthogonal to it. 
A target pallet capable of containing \(n\) boxes is given.
Each box \(i\) has a destination position on the pallet, given by 
$(P_x[i], P_y[i])$, where $i = 1,\ldots,n$. 

The robot can pick a consecutive set of boxes from the stream. Since boxes
on the belt have no gaps between them, each picked set is contiguous
along the \(x\)-axis. Given a gripper size \(S\), the total length of the
picked set must be at most \(S+x\). This corresponds to an overflow of at
most \(x/2\) on each side of the gripper. Therefore, the maximum number
of boxes that can be picked at once is $\left\lfloor \frac{S+x}{x} \right\rfloor$.

After picking a set of boxes, the robot releases it on the pallet.
Figure~\ref{fig:picking_releasing} shows an example of the picking
and releasing operations.

\begin{figure}[h]
    \centering
    \includegraphics[width=0.75\textwidth]{figures/pick_and_release.pdf}
    \caption{Example of picking and releasing operations.}
    \label{fig:picking_releasing}
\end{figure}
The gripper has two parallel plates, but only one plate can move during
release. For this reason, the opening side must be free. It can face either
the outside of the pallet or a position that has not been filled yet.
We split the problem into two parts:
\begin{enumerate}
    \item \textbf{Grouping.} Identify the sets of boxes that can be picked
    together. Each group must contain boxes on the same row, aligned along
    the \(x\)-axis, and must fit within the gripper capacity.

    \item \textbf{Scheduling.} Choose the order in which the groups are
    released. For each group, also choose which plate opens. At every step,
    the selected opening side must be free.
\end{enumerate}
The goal is to find a sequence of pick-and-release actions that places all
boxes in their target positions while respecting the gripper capacity
and avoiding collisions during release. The code used for the model
and the benchmark instances is available at \url{...}.

\subsection{Parameters and Derived Constants}

The proposed MiniZinc model uses five input parameters. The parameter
\texttt{n\_boxes} denotes the total number of boxes in the layer.
The parameters \texttt{boxes\_along\_x} and \texttt{boxes\_along\_y}
define the grid dimensions of the layer. The parameter
\texttt{boxes\_x\_dim} denotes the box dimension parallel to the gripper.
In the natural orientation, this corresponds to \(x\); if the layer is
rotated, it corresponds to \(y\). Finally, \texttt{grip\_size} denotes
the gripper size \(S\). These parameters are declared as follows:
\vspace{0.5cm}
\begin{lstlisting}
par int: n_boxes;
par int: boxes_x_dim;
par int: boxes_along_x;
par int: boxes_along_y;
par int: grip_size;
\end{lstlisting}
\vspace{0.5cm}
From these inputs, the model derives three constants. The first one is the
maximum number of boxes that can be picked in a single operation, computed as:
\[
    \texttt{max\_pkble\_boxes}
    =
    \left\lfloor
    \frac{\texttt{grip\_size} + \texttt{boxes\_x\_dim}}
         {\texttt{boxes\_x\_dim}}
    \right\rfloor .
\]
This value derives directly from the gripper capacity \(S+x\). The second
constant is the number of groups needed to cover one pallet row:
\[
    \texttt{groups\_per\_row}
    =
    \left\lceil
    \frac{\texttt{boxes\_along\_x}}
         {\texttt{max\_pkble\_boxes}}
    \right\rceil .
\]
This is a ceiling division because the last group in a row may contain fewer
boxes than the maximum allowed number. The third constant is the total number of groups:
\[
    \texttt{n\_groups}
    =
    \texttt{boxes\_along\_y}
    \cdot
    \texttt{groups\_per\_row}.
\]
It is obtained by multiplying the number of groups per row by the number of
rows. The corresponding MiniZinc declarations are:
\vspace{0.5cm}
\begin{lstlisting}
par int: max_pkble_boxes = (grip_size + boxes_x_dim) div boxes_x_dim;

par int: groups_per_row = (boxes_along_x + max_pkble_boxes - 1)
                           div max_pkble_boxes;

par int: n_groups = boxes_along_y * groups_per_row;
\end{lstlisting}